%这是一个demo摘要，请替换成你的摘要。

本文针对北京工业大学耿丹学院本科生毕业设计论文的排版需求，设计并实现了一套基于 LaTeX 的自动化排版系统。系统严格按照学校毕业设计模板的格式要求，对页面设置、字体字号、章节标题、图表公式、参考文献等进行了精细化配置，形成了可直接使用的文档类模板。

研究工作首先分析了学校 Word 模板的排版规范，包括各级标题的字号与行距、正文的缩进与段间距、图表的编号规则以及参考文献的著录格式。在此基础上，利用 ctex 宏包实现了中文字体的配置，利用 titlesec 和 titletoc 宏包实现了标题与目录格式的定制，利用 gbt7714 宏包实现了参考文献的国标格式。最终，系统通过一个文档类文件（.cls）将所有配置整合，用户仅需关注内容撰写即可自动生成符合规范的论文。

实验结果表明，该模板能够准确复现学校模板的排版效果，且具备良好的稳定性与可移植性，可显著提升毕业设计论文的排版效率与规范性。
